\documentclass[10pt,a4paper]{amsart}
\usepackage[margin=19mm]{geometry}
\usepackage{amsmath,amssymb,mathtools}
\usepackage[scr=rsfs,cal=euler]{mathalfa}
\usepackage{xfrac}
\usepackage[unicode,hidelinks,bookmarks=false]{hyperref}
\newcommand{\ie}{i.e.\ }
\newcommand{\cf}{cf.\ }
\newcommand{\resp}{respectively\ }
\newcommand{\Subsection}{\subsection}
\providecommand{\nobreakdash}{\nobreak\hbox{-}}
\setlength{\emergencystretch}{2em}
\multlinegap0pt
\usepackage{cancel}
\usepackage{amssymb}
\usepackage{xcolor}
\newtheorem{lem}{Lemma}
\newtheorem{prop}{Proposition}
\title{Perverse sheaves on the semi-infinite flag variety and representations of the Frobenius kernel}
\author{Emilien Zabeth}
\date{Source: arXiv:2504.03341v2 (CC BY 4.0)}
\begin{document}
\maketitle

\noindent\emph{Source and licence.} Perverse sheaves on the semi-infinite flag variety and representations of the Frobenius kernel. Version arXiv:2504.03341v2 (CC BY 4.0), distributed by arXiv under the Creative Commons Attribution 4.0 International licence, http://creativecommons.org/licenses/by/4.0/, as recorded in the arXiv licence metadata for this exact version and shown on the abstract page https://arxiv.org/abs/2504.03341v2. Full licence notice: \texttt{licenses/arxiv-2504.03341-CC-BY-4.0.txt}.\par\smallskip

\noindent\textit{Editorial scope.} Counted unit: the article's lem convo (source0.tex lines 1763--1766). The article's complete original proof of it is delivered with it (source0.tex lines 1767--1779, one \texttt{proof} environment of 13 lines). No mathematics is written, repaired or re-derived: the statement and the proof are the article's own lines, in the article's own notation, and no retained line is substituted or re-flowed.  Retained uncounted as the source-local context the proof needs: lines 356--363; lines 652--655; lines 667--670; lines 1643--1645.  Nothing was omitted from any retained span, so the package carries no omission record.  No retained line contains a citation, so the wrapper carries no bibliography and no bibliography entry.  No retained line contains a figure environment, an image-inclusion command or drawing code, so no image is delivered and the package declares no asset dependency.  Editorial formatting only, in addition: an AMS \texttt{amsart} preamble that restates the article's own declarations and macros used by the retained lines, namely the environments \texttt{lem}, \texttt{prop} on separate counters, together with the standard packages those lines need.  The retained environments print in this document as Lemma 1, Proposition 1 in order of appearance, whereas the article prints the same items with the numbering of its own full document.  The complete native source of the article as distributed in the arXiv e-print for this exact version is bundled unchanged as \texttt{sources/175965/source0.tex}.
	\begin{lem}\label{lem conv standard}
 \begin{enumerate}
     \item For any $w\in W_\mathrm{ext}$, there exist canonical isomorphisms
     $$\mathcal{D}_w\star^I\mathcal{N}_{w^{-1}}\simeq \mathcal{D}_e\simeq \mathcal{N}_w\star^I\mathcal{D}_{w^{-1}}. $$
     \item For $w,y\in W_\mathrm{ext}$ such that $\ell(ww_0y)=\ell(w)+\ell(w_0y)$ and both $wy$ and $y$ belong to $W_\mathrm{ext}^S$, there exist canonical isomorphisms
		$$\mathcal{D}_w\star^I\Delta_y\simeq\Delta_{wy}\qquad \mathcal{N}_w\star^I\nabla_y\simeq\nabla_{wy}.$$
 \end{enumerate}
	\end{lem}

	\begin{prop}\label{prop convolution costandard}
		For any  $\gamma\in\mathbf{Y}^+,~\nu\in\mathbf{Y}$, we have a canocical isomorphism
		$$ \mathcal{N}_\gamma\star^I\widehat{\mathcal{Z}}'_\nu\simeq \widehat{\mathcal{Z}}'_{\gamma+\nu} $$
	\end{prop}

 \begin{prop}\label{prop convolution costandard 2}
     For any $w,w'\in W$ such that $\ell(ww'w_0)=\ell(w)+\ell(w'w_0)$ and $\lambda\in \mathbf{Y}$, we have a canonical isomorphism
     $$ \mathcal{N}_{w}\star^I\widehat{\mathcal{Z}}'_{w't_\lambda}\simeq \widehat{\mathcal{Z}}'_{ww't_\lambda}. $$
 \end{prop}

	\begin{equation}\label{diagram coeq}
		\mathcal{S}_\lambda\star\mathcal{R}_{-\mu}\star\mathcal{IC}^{\frac{\infty}{2}}_{\mu-\lambda}\rightrightarrows\mathcal{F}, 
	\end{equation}

\begin{lem}\label{lem commute convo}
    Let $\gamma\in\mathbf{Y}^+$, $\nu\in\mathbf{Y}$ and  $w,w'\in W$ such that $\ell(ww_0w')=\ell(w)+\ell(w_0w')$. We have canonical isomorphisms
    $$ \mathcal{N}_\gamma\star^I\mathrm{Conv}^{\mathrm{Hecke}}(\widehat{\mathcal{Z}}'_\nu)\simeq \mathrm{Conv}^{\mathrm{Hecke}}(\widehat{\mathcal{Z}}'_{\gamma+\nu}),\quad\mathcal{N}_w\star^I\mathrm{Conv}^{\mathrm{Hecke}}(\widehat{\mathcal{Z}}'_{w't_\nu})\simeq \mathrm{Conv}^{\mathrm{Hecke}}(\widehat{\mathcal{Z}}'_{ww't_\nu}).$$
\end{lem}

\begin{proof}
   Recall that  $\mathrm{Conv}^{\mathrm{Hecke}}(\widehat{\mathcal{Z}}'_{\gamma+\nu})$ represents the covariant functor which, to any $\mathcal{F}\in \mathrm{Ind}(\mathrm{Perv}_{{I_\mathrm{u}}}(\mathcal{F}l^{\frac{\infty}{2}}))$ , associates the set of collections of morphisms $$((\widehat{\mathcal{Z}}'_{\gamma+\nu})_\lambda\star\mathcal{IC}^{\frac{\infty}{2}}_{-\lambda}\to \mathcal{F})_{\nu\in\mathbf{Y}}$$ such that, for any $\lambda,\mu\in\mathbf{Y}$, the following diagram commutes (cf. \eqref{diagram coeq}):
	\begin{equation*}
		(\widehat{\mathcal{Z}}'_{\gamma+\nu})_\lambda\star\mathcal{R}_{-\mu}\star\mathcal{IC}^{\frac{\infty}{2}}_{\mu-\lambda}\rightrightarrows\mathcal{F}. 
	\end{equation*} 
 Applying Proposition \ref{prop convolution costandard} together with the first point of Lemma \ref{lem conv standard}, we obtain that $\mathrm{Conv}^{\mathrm{Hecke}}(\widehat{\mathcal{Z}}'_{\gamma+\nu})$ represents the covariant functor which, to any $\mathcal{F}\in \mathrm{Ind}(\mathrm{Perv}_{{I_\mathrm{u}}}(\mathcal{F}l^{\frac{\infty}{2}}))$ , associates the set of collections of morphisms $$((\widehat{\mathcal{Z}}'_{\nu})_\lambda\star\mathcal{IC}^{\frac{\infty}{2}}_{-\lambda}\to \mathcal{D}_{-\gamma}\star\mathcal{F})_{\nu\in\mathbf{Y}}$$ such that, for any $\lambda,\mu\in\mathbf{Y}$, the following diagram commutes (cf. \eqref{diagram coeq}):
	\begin{equation*}
		(\widehat{\mathcal{Z}}'_{\nu})_\lambda\star\mathcal{R}_{-\mu}\star\mathcal{IC}^{\frac{\infty}{2}}_{\mu-\lambda}\rightrightarrows\mathcal{D}_{-\gamma}\star\mathcal{F}. 
	\end{equation*} 
 By the first point of Lemma \ref{lem conv standard} again, this functor is represented by $ \mathcal{N}_\gamma\star^I\mathrm{Conv}^{\mathrm{Hecke}}(\widehat{\mathcal{Z}}'_\nu)$.

  The second isomorphism is proved in the same way, using that $\mathcal{N}_w\star^I\widehat{\mathcal{Z}}'_{w't_\nu}\simeq \widehat{\mathcal{Z}}'_{ww't_\nu}$ (Proposition \ref{prop convolution costandard 2}). 
\end{proof}

\end{document}
